\documentclass{article}
\begin{document}

\section{Abstract}
We study attention-temperature self-distillation for LLM unlearning.
We find no defensible improvement from spatial temperature fields, temporal
temperature schedules, or forget-gradient subspace projections once the mean
temperature budget is held fixed. A non-equal-budget comparison had suggested
a 0.1493 NLL gain for a low2high schedule over global tau=2, but the equal-budget
control at tau=2.25 reaches dF=1.1416 while low2high reaches dF=0.9491.
Our contribution is therefore an evaluation protocol: an equal-budget control
plus a P1/P2/P3/INT8 recovery ladder showing that 0.4477 of the forgetting gain
can be undone by few-shot prompting or 20 steps of SGD relearning.

\section{Method}
The agent may edit only \texttt{train.py}; \texttt{prepare.py} and the fixed
evaluation function are read-only. Round 1 uses 8 experiments with the
equal-budget shrink regret as the single metric. Round 2 uses 6 experiments
with the depth regret as the single metric. All rounds use 60 steps,
alpha=0.1, learning rate 5e-6, seed 0, on TOFU forget10 and retain90 with a
Llama-3.2-1B-Instruct base model. Peak memory is 51.1 GB in round 1 and
51.3 GB in round 2.

\section{Results}
Table~\ref{tab:r1} reports all eight round-1 experiments and Table~\ref{tab:r2}
reports all six round-2 experiments. Zero candidates were kept in round 1 and
zero candidates were kept in round 2. The baseline val_bpb is 0.000000 and the
best val_bpb is 0.000000 because both are the anchor round itself.

\begin{table}[h]
\caption{Round 1 experiment log (equal-budget shrink regret)}
\label{tab:r1}
\begin{tabular}{lrrr}
round & val\_bpb & memory\_gb & status \\
A1 & 0.000000 & 51.1 & keep \\
A2 & 0.000000 & 51.1 & keep \\
C1 & 0.302613 & 51.1 & discard \\
C2 & 0.599510 & 51.1 & discard \\
C3 & 0.251062 & 51.1 & discard \\
C4 & 0.816921 & 51.1 & discard \\
C5 & 0.093430 & 51.1 & discard \\
C6 & 0.314867 & 51.1 & discard \\
\end{tabular}
\end{table}

\begin{table}[h]
\caption{Round 2 experiment log (depth regret)}
\label{tab:r2}
\begin{tabular}{lrrr}
round & val\_bpb & memory\_gb & status \\
B0 & 0.000000 & 51.3 & keep \\
B1 & 0.027804 & 51.3 & discard \\
B2 & 0.022247 & 51.3 & discard \\
B3 & 0.028512 & 51.3 & discard \\
B4 & 0.027469 & 51.3 & discard \\
B5 & 0.029642 & 51.3 & discard \\
\end{tabular}
\end{table}

\section{Matched-dF control}
The hinge anti-recovery candidate B5 reaches rho=0.408529 with dF=1.8046, while
a plain global tau=2.607 control B6 reaches rho=0.590908 with dF=1.7914.
The difference of 0.182379 is not explained by denominator inflation.
However B5 was still discarded by the fixed metric because the complexity
penalty is 0.040000 and the third term of the metric is asymmetric.

\section{Limitations}
We use a single seed, 60 steps, one 1B model on TOFU forget10, and one P3 attack
strength with one optimizer. The hinge candidate may produce degenerate outputs
rather than erasing knowledge; generation quality, refusal rate and 3-seed
repeats were not measured. The dual-loop meta objective of R2-C is NOT implemented;
B5 is only a single-loop hinge approximation. No independent reproduction and no
cross-model reviewer were available. Projection operators were computed once
before training rather than per step.

\end{document}
